El model de l'àtom d'hidrogen de Bohr suposa que l'electró descriu una òrbita circular al voltant del protó i que el moment angular de l'òrbita és un múltiple enter de la constant de Planck $L = n\cdot\dfrac{h}{2\pi}$, on $n$ és un nombre natural.

\begin{apartat}{1,25}
Demostreu que el moment angular s'expressa com $L = e\sqrt{k\cdot m\cdot r}$, on $e$ és la càrrega de l'electró, $k$ és la constant de Coulomb, $m$ és la massa de l'electró i $r$ és el radi de l'òrbita. Trobeu l'expressió del radi de l'òrbita de l'electró en funció de $h$, $n$, $k$, $m$ i $e$. Calculeu el radi per a $n = 1$.
\end{apartat}

\begin{apartat}{1,25}
Suposant que el radi de l'òrbita per a $n = 1$ és 53~pm, calculeu les energies cinètica, potencial i mecànica clàssiques de l'electró.
\end{apartat}

\dades{%
  $|e| = 1,602\times 10^{-19}\ \mathrm{C}$.\\
  $k = \dfrac{1}{4\pi\varepsilon_0} = 8,99\times 10^{9}\ \mathrm{N\,m^2\,C^{-2}}$.\\
  $m = 9,11\times 10^{-31}\ \mathrm{kg}$.\\
  $h = 6,63\times 10^{-34}\ \mathrm{J\,s}$.}
